\begin{figure}[htbp]
			\centering
			\resizebox{\textwidth}{!}{%
				\begin{tikzpicture}[
					scale=0.95,
					>=latex,
					trave/.style={line width=1.2pt},
					carico/.style={black, line width=0.4pt},
					forza/.style={black, line width=1pt},
					risultante/.style={gray!65, line width=1.8pt},
					diagT/.style={black, line width=0.8pt},
					diagM/.style={black, line width=0.8pt},
					diagTb/.style={gray!65, line width=1.0pt},
					diagMb/.style={gray!65, line width=1.0pt},
					hatch/.style={gray!40, line width=0.3pt},
					quota/.style={black, line width=0.4pt},
					fibre/.style={gray!70, line width=0.3pt},
					segno/.style={circle, draw=black, line width=0.4pt,
						minimum size=10pt, inner sep=0pt, font=\small\bfseries},
					]
					
					% ── Parametri dei vincoli (convenzione C06/C07) ──
					\def\tw{0.15}    \def\thR{0.28}   \def\rn{0.12}
					\def\rw{0.05}    \def\bw{0.09}    \def\gw{0.22}
					\def\hh{0.10}
					
					% ── Geometria del problema ──
					\def\Lell{3}                            % ℓ = lunghezza campata BC
					\def\Lsb{1.5}                           % ℓ/2 = lunghezza di ciascuno sbalzo
					\pgfmathsetmacro{\xA}{0}                % estremo libero sinistro
					\pgfmathsetmacro{\xB}{\Lsb}             % cerniera interna
					\pgfmathsetmacro{\xMid}{\xB + 0.5*\Lell}% mezzeria BC = ℓ
					\pgfmathsetmacro{\xC}{\xB + \Lell}      % carrello interno
					\pgfmathsetmacro{\xD}{\xC + \Lsb}       % estremo libero destro
					\pgfmathsetmacro{\xMidAB}{0.5*\Lsb}     % mezzeria AB
					\pgfmathsetmacro{\xMidBmid}{0.5*(\xB+\xMid)} % mezzo di B-mezzeria
					\pgfmathsetmacro{\xMidmidC}{0.5*(\xMid+\xC)} % mezzo di mezzeria-C
					\pgfmathsetmacro{\xMidCD}{0.5*(\xC+\xD)}    % mezzeria CD
					\def\yqi{-0.80}
					\def\ylab{-2.80}
					
					% ── Macro vincoli ──
					\def\cerniera#1{%
						\draw[thick] (#1,0) -- (#1-\tw,-\thR) -- (#1+\tw,-\thR) -- cycle;
						\draw[thick] (#1-\gw,-\thR) -- (#1+\gw,-\thR);
						\foreach \xh in {-0.16,-0.08,0,0.08,0.16}{%
							\draw[thin] (#1+\xh,-\thR) -- ++(-0.07,-\hh);}
						\draw[thick, fill=white] (#1,0) circle (\rn);}
					% variante per cerniera in punto INTERNO della trave
					\def\cernieraInt#1{%
						\draw[thick] (#1,0) -- (#1-\tw,-\thR) -- (#1+\tw,-\thR) -- cycle;
						\draw[thick] (#1-\gw,-\thR) -- (#1+\gw,-\thR);
						\foreach \xh in {-0.16,-0.08,0,0.08,0.16}{%
							\draw[thin] (#1+\xh,-\thR) -- ++(-0.07,-\hh);}
						\draw[thick, fill=white] (#1-\rn,0) arc(180:360:\rn) -- cycle;}
					% variante per carrello in punto INTERNO della trave
					\def\carrelloInt#1{%
						\draw[thick] (#1,0) -- (#1-\tw,-\thR) -- (#1+\tw,-\thR) -- cycle;
						\draw[thick] (#1-\bw,-\thR-\rw-0.01) circle (\rw);
						\draw[thick] (#1+\bw,-\thR-\rw-0.01) circle (\rw);
						\draw[thick] (#1-\gw,-\thR-2*\rw-0.04) -- (#1+\gw,-\thR-2*\rw-0.04);
						\foreach \xh in {-0.16,-0.08,0,0.08,0.16}{%
							\draw[thin] (#1+\xh,-\thR-2*\rw-0.04) -- ++(-0.07,-\hh);}
						\draw[thick, fill=white] (#1-\rn,0) arc(180:360:\rn) -- cycle;}
					
					% ── Carico distribuito (riquadro) ──
					\def\pgap{0.10}    \def\php{0.40}
					\pgfmathsetmacro{\pBase}{\pgap}
					\pgfmathsetmacro{\pTop}{\pgap+\php}
					
					% ── Fattori di scala ──
					\def\sT{1.8}
					\def\sM{7.0}
					
					% =========================================================
					% PANNELLO (a) — Schema strutturale
					% =========================================================
					\begin{scope}[xshift=0cm]
						
						\draw[trave] (\xA,0) -- (\xD,0);
						\draw[fibre, dashed] (\xA, -0.1) -- (\xD, -0.1);
						
						\cernieraInt{\xB}
						\carrelloInt{\xC}
						
						% Riquadro del carico distribuito p̄ su tutta la trave
						\draw[carico] (\xA,\pBase) rectangle (\xD,\pTop);
						\foreach \x in {0.25, 0.75, 1.25, 1.75, 2.25, 2.75, 3.25, 3.75,
							4.25, 4.75, 5.25, 5.75} {
							\draw[carico, -{Latex[length=2pt, width=1.5pt]}]
							(\x, \pTop) -- (\x, \pBase);
						}
						\node[above, font=\small] at (0.5, \pTop+0.02) {$\bar{p}$};
						
						% Risultanti grigie (sostituzione caso b: ogni distribuito → propria risultante)
						\def\Rstart{1.5}
						% Risultante sullo sbalzo AB
						\draw[risultante, -{Latex[length=5pt, width=4pt]}]
						(\xMidAB, \Rstart) -- (\xMidAB, 0);
						\node[left, font=\small, gray!50!black] at (\xMidAB-0.06, \Rstart-0.18)
						{$\bar p\ell/2$};
						% Risultante sulla campata BC
						\draw[risultante, -{Latex[length=5pt, width=4pt]}]
						(\xMid, \Rstart) -- (\xMid, 0);
						\node[right, font=\small, gray!50!black] at (\xMid+0.08, \Rstart-0.18)
						{$\bar p\ell$};
						% Risultante sullo sbalzo CD
						\draw[risultante, -{Latex[length=5pt, width=4pt]}]
						(\xMidCD, \Rstart) -- (\xMidCD, 0);
						\node[right, font=\small, gray!50!black] at (\xMidCD+0.08, \Rstart-0.18)
						{$\bar p\ell/2$};
						
						\node[above left=-2pt,  font=\small] at (\xA,0)   {$A$};
						\node[above, font=\small] at (\xB, 0.04)          {$B$};
						\node[above, font=\small] at (\xC, 0.04)          {$C$};
						\node[above right=-2pt, font=\small] at (\xD,0)   {$D$};
						
						% Quote orizzontali
						\draw[<->, quota] (\xA,\yqi) -- (\xB,\yqi);
						\draw[<->, quota] (\xB,\yqi) -- (\xC,\yqi);
						\draw[<->, quota] (\xC,\yqi) -- (\xD,\yqi);
						\draw[quota] (\xA,\yqi+0.07) -- (\xA,\yqi-0.07);
						\draw[quota] (\xB,\yqi+0.07) -- (\xB,\yqi-0.07);
						\draw[quota] (\xC,\yqi+0.07) -- (\xC,\yqi-0.07);
						\draw[quota] (\xD,\yqi+0.07) -- (\xD,\yqi-0.07);
						\node[above, font=\footnotesize] at (\xMidAB, \yqi-0.08) {$\ell/2$};
						\node[above, font=\footnotesize] at (\xMid, \yqi-0.08) {$\ell$};
						\node[above, font=\footnotesize] at (\xMidCD, \yqi-0.08) {$\ell/2$};
						
						\node[font=\small] at (\xD/2, \ylab/2-0.25) {(a)};
						
					\end{scope}
					
					% =========================================================
					% PANNELLO (b) — Diagramma del taglio T(z)
					% =========================================================
					\begin{scope}[xshift=7.5cm]
						
						\draw[trave] (\xA,0) -- (\xD,0);
						\cernieraInt{\xB}
						\carrelloInt{\xC}
						
						\node[above left=-2pt,  font=\small] at (\xA,0)   {$A$};
						\node[above left, font=\small] at (\xB, 0.04)          {$B$};
						\node[above left, font=\small] at (\xC, 0.04)          {$C$};
						\node[above right=-2pt, font=\small] at (\xD,0)   {$D$};
						
						\pgfmathsetmacro{\Tval}{0.5*\sT}    % |T_max| = p̄ℓ/2 → 0.8
						
						% Hatching del caso (a) — lineare a tratti
						\foreach \x in {0.10, 0.30, ..., 1.40} {
							\pgfmathsetmacro{\h}{-\sT*\x/\Lell}
							\draw[hatch] (\x, 0) -- (\x, \h);
						}
						\foreach \x in {1.60, 1.80, ..., 4.40} {
							\pgfmathsetmacro{\h}{\sT*(\xMid - \x)/\Lell}
							\draw[hatch] (\x, 0) -- (\x, \h);
						}
						\foreach \x in {4.60, 4.80, ..., 5.90} {
							\pgfmathsetmacro{\h}{\sT*(\xD - \x)/\Lell}
							\draw[hatch] (\x, 0) -- (\x, \h);
						}
						
						% Caso (b) GRIGIO: spezzata costante a tratti — sostituzione completa
						% Salti in mezzeria AB, in B, in mezzeria BC, in C, in mezzeria CD
						\draw[diagTb]
						(\xA, 0) -- (\xMidAB, 0)
						-- (\xMidAB, -\Tval) -- (\xB, -\Tval)
						-- (\xB, \Tval) -- (\xMid, \Tval)
						-- (\xMid, -\Tval) -- (\xC, -\Tval)
						-- (\xC, \Tval) -- (\xMidCD, \Tval)
						-- (\xMidCD, 0) -- (\xD, 0);
						
						% Caso (a) NERO: lineare a tratti continuo
						\draw[diagT]
						(\xA, 0) -- (\xB, -\Tval)
						-- (\xB, \Tval) -- (\xC, -\Tval)
						-- (\xC, \Tval) -- (\xD, 0);
						
						% Segni delle aree
						\node[segno] at (\xMidAB, {-0.5*\Tval}) {$-$};
						\node[segno] at (\xMidBmid, {0.5*\Tval}) {$+$};
						\node[segno] at (\xMidmidC, {-0.5*\Tval}) {$-$};
						\node[segno] at (\xMidCD, {0.5*\Tval}) {$+$};
						
						% Etichette valori notevoli
						\node[left, font=\small] at (\xB-0.02, \Tval-0.02){$\tfrac{\bar p\ell}{2}$};
						\node[right, font=\small] at (\xB-0.02, -\Tval+0.02){$\tfrac{\bar p\ell}{2}$};
						\node[right, font=\small] at (\xC+0.02, -\Tval+0.02){$\tfrac{\bar p\ell}{2}$};
						\node[left, font=\small] at (\xC+0.02, \Tval+0.02){$\tfrac{\bar p\ell}{2}$};
						\node[below, font=\footnotesize] at (\xMid, -0.05) {$z_0$};
						
						\node[font=\small] at (\xD/2, \ylab/2-0.25) {(b)};
						
					\end{scope}
					
					% =========================================================
					% PANNELLO (c) — Diagramma del momento M(z)
					% =========================================================
					\begin{scope}[xshift=15.0cm]
						
						\draw[trave] (\xA,0) -- (\xD,0);
						\cernieraInt{\xB}
						\carrelloInt{\xC}
						
						\node[above left=-2pt,  font=\small] at (\xA,0)   {$A$};
						\node[above, font=\small] at (\xB, 0.04)          {$B$};
						\node[above, font=\small] at (\xC, 0.04)          {$C$};
						\node[above right=-2pt, font=\small] at (\xD,0)   {$D$};
						
						\pgfmathsetmacro{\Mval}{\sM/8}     % |M_max| = p̄ℓ²/8 → 0.625
						
						% Hatching del caso (a) — parabole nei tre tratti
						\foreach \x in {0.05, 0.15, ..., 1.45} {
							\pgfmathsetmacro{\h}{(\sM/(2*\Lell*\Lell))*\x*\x}
							\draw[hatch] (\x, 0) -- (\x, \h);
						}
						% BC
						\foreach \x in {1.55, 1.65, ..., 4.45} {
							\pgfmathsetmacro{\h}{(\sM/(2*\Lell*\Lell))*(\x - \xMid)*(\x - \xMid)}
							\draw[hatch] (\x, 0) -- (\x, \h);
						}
						% CD
						\foreach \x in {4.55, 4.65, ..., 5.95} {
							\pgfmathsetmacro{\h}{(\sM/(2*\Lell*\Lell))*(\xD - \x)*(\xD - \x)}
							\draw[hatch] (\x, 0) -- (\x, \h);
						}
						
						% Caso (b) GRIGIO: spezzata lineare a tratti su TUTTA la trave
						% Vertici nei punti notevoli A, mezz.AB, B, mezz.BC, C, mezz.CD, D
						% Ogni segmento è tangente alla parabola caso (a) in un punto notevole
						\draw[diagMb]
						(\xA, 0) -- (\xMidAB, 0) -- (\xB, \Mval) -- (\xMid, -\Mval)
						-- (\xC, \Mval) -- (\xMidCD, 0) -- (\xD, 0);
						
						% Caso (a) NERO: tre parabole concave verso l'alto
						% AB: parte da A in (0,0) con tang. orizz., arriva in B a (xB, Mval)
						\draw[diagM]
						(\xA, 0)
						\foreach \x in {0.10, 0.20, ..., 1.40} {
							-- ({\x}, {(\sM/(2*\Lell*\Lell))*\x*\x})
						}
						-- (\xB, \Mval);
						% BC: parte da B a (xB, Mval), tocca asse in mezzeria con tang. orizz.,
						% risale a C a (xC, Mval)
						\draw[diagM]
						(\xB, \Mval)
						\foreach \x in {1.60, 1.70, ..., 4.40} {
							-- ({\x}, {(\sM/(2*\Lell*\Lell))*(\x - \xMid)*(\x - \xMid)})
						}
						-- (\xC, \Mval);
						% CD: parte da C a (xC, Mval), arriva in D a (xD, 0) con tang. orizz.
						\draw[diagM]
						(\xC, \Mval)
						\foreach \x in {4.60, 4.70, ..., 5.90} {
							-- ({\x}, {(\sM/(2*\Lell*\Lell))*(\xD - \x)*(\xD - \x)})
						}
						-- (\xD, 0);
						
						% Corda dei momenti d'estremità della campata centrale BC
						\draw[diagMb, dashed] (\xB, \Mval) -- (\xC, \Mval);
						
						
						% Etichette valori notevoli (M negativo → segno meno; M positivo → senza segno)
						\node[left, font=\small] at (\xB-0.02, \Mval-0.02){$\tfrac{\bar p\ell^2}{8}$};
						\node[right, font=\small] at (\xC+0.02, \Mval-0.02){$\tfrac{\bar p\ell^2}{8}$};
						\node[right, font=\footnotesize, gray!50!black]	at (\xMid, -\Mval-0.02) {$\tfrac{\bar p\ell^2}{8}$};
						
						\node[font=\small] at (\xD/2, \ylab/2-0.25) {(c)};
						
					\end{scope}
					
				\end{tikzpicture}
			}
			\caption{Trave con doppio sbalzo simmetrico e carico uniforme $\bar p$
				(caso~a, tratto nero) confrontata con la sostituzione di ogni
				carico distribuito con la propria risultante applicata nel
				baricentro del tratto: $\bar p\ell/2$ in mezzeria di $AB$ e di
				$CD$, $\bar p\ell$ in mezzeria di $BC$ (caso~b, tratto grigio).
				(a)~schema strutturale; (b)~diagrammi del taglio; (c)~diagrammi
				del momento. Ogni segmento della spezzata grigia del momento è
				tangente alla parabola caso~(a) in uno dei punti
				$\{A,\,B^-,\,B^+,\,C^-,\,C^+,\,D\}$. Sul tratto centrale $BC$
				è inoltre tracciata (tratto grigio tratteggiato) la corda dei
				momenti d'estremità, retta orizzontale a quota $-\bar p\ell^2/8$;
				la tangente parallela alla parabola in mezzeria di $BC$ coincide
				con l'asse della trave (poiché $T(\ell) = 0$) e non è disegnata
				separatamente.}
			\label{fig:stat_sbalzo_02}
		\end{figure}
